\documentclass[aspectratio=169, lualatex, handout]{beamer}
\makeatletter\def\input@path{{theme/}}\makeatother\usetheme{cipher}

\title{Applied Cryptography}
\author{Nadim Kobeissi}
\institute{American University of Beirut}
\instituteimage{images/aub_white.png}
\date{\today}
\coversubtitle{CMPS 297AD/396AI}
\coverpartname{Part 1: Provable Security}
\coversessionname{1.7: Hard Problems}
\coverwebsite{https://appliedcryptography.page}

\begin{document}
\begin{frame}[plain]
	\titlepage
\end{frame}

\begin{frame}{How it's made}
	\bigimagewithcaption{fischer.png}{Fischer et al., The Challenges of Bringing Cryptography from Research Papers to Products: Results from an Interview Study with Experts, USENIX Security 2024}
\end{frame}

\begin{frame}{Cryptographic building blocks}
	\begin{columns}[c]
		\begin{column}{0.5\textwidth}
			\textbf{Security goals}
			\begin{itemize}[<+->]
				\item \textbf{Confidentiality}: Data exchanged between Client and Server
				      is only known to those parties.
				\item \textbf{Authentication}: If Server receives data from Client,
				      then Client sent it to Server.
				\item \textbf{Integrity}: If Server modifies data owned by Client,
				      Client can find out.
			\end{itemize}
		\end{column}

		\begin{column}{0.5\textwidth}
			\textbf{Examples}
			\begin{itemize}[<+->]
				\item \textbf{Confidentiality}: When you send a private message on Signal,
				      only you and the recipient can read the content.
				\item \textbf{Authentication}: When you receive an email from your boss,
				      you can verify it actually came from them.
				\item \textbf{Integrity}: Your computer can verify that software update
				      downloads haven't been tampered with during transmission.
			\end{itemize}
		\end{column}
	\end{columns}
\end{frame}

\begin{frame}{Security goals: more examples}
	\begin{itemize}[<+->]
		\item \textbf{TLS (HTTPS)} ensures that data exchanged between the client
		      and the server is confidential and that parties are authenticated.
		      \begin{itemize}
			      \item Allows you to log into gmail.com without your ISP learning your password.
		      \end{itemize}
		\item \textbf{FileVault 2} ensures data confidentiality and integrity on
		      your MacBook.
		      \begin{itemize}
			      \item Prevents thieves from accessing your data if your MacBook is stolen.
		      \end{itemize}
		\item \textbf{Signal} implements post-compromise security, an advanced security
		      goal.
		      \begin{itemize}
			      \item Allows a conversation to ``heal'' in the event of a temporary key
			            compromise.
			      \item More on that later in the course.
		      \end{itemize}
	\end{itemize}
\end{frame}

\begin{frame}{Why bother?}
	\begin{itemize}[<+->]
		\item Can't we just use access control?
		\item Strictly speaking, usernames and passwords can be implemented
		      without cryptography\ldots
		\item Server checks if the password matches, or if the IP address matches,
		      etc. before granting access.
		\item What's so bad about that?
	\end{itemize}
	\definitionbox{The Problem with Traditional Access Control}{
		\begin{itemize}[<+->]
			\item Requires trusting the server completely
			\item No protection during transmission
			\item No way to verify integrity
			\item No way to establish trust between strangers
		\end{itemize}
	}
\end{frame}

\begin{frame}[c]{The magic of cryptography}
	\begin{center}
		\Large\textbf{Cryptography lets us achieve what seems impossible}
		\vspace{1cm}
		\begin{itemize}[<+->]
			\item Secure communication over insecure channels
			\item Verification without revealing secrets
			\item Proof of computation without redoing it
		\end{itemize}
	\end{center}
\end{frame}

\begin{frame}{Hard problems}
	\begin{itemize}[<+->]
		\item Cryptography is largely about equating the security of a system to the
		      difficulty of solving a math problem that is thought to be computationally
		      very expensive.
		\item With cryptography, we get security systems that we can literally
		      mathematically prove as secure (under assumptions).
		\item Also, this allows for actual magic.
		      \begin{itemize}[<+->]
			      \item Alice and Bob meet for the first time in the same room as you.
			      \item You are listening to everything they are saying.
			      \item Can they exchange a secret without you learning it?
		      \end{itemize}
	\end{itemize}
\end{frame}

\begin{frame}{Time for actual magic}
	\bigimagewithcaption{dh.png}{}
\end{frame}

\begin{frame}{No known feasible computation}
	\begin{itemize}[<+->]
		\item The discrete logarithm problem:
		      \begin{itemize}
			      \item Given a finite cyclic group $G$, a generator $g \in G$, and an element
			            $h \in G$, find the integer $x$ such that $g^{x}=h$
		      \end{itemize}
		\item In more concrete terms:
		      \begin{itemize}
			      \item Let $p$ be a large prime and let $g$ be a generator of the multiplicative
			            group $\mathbb{Z}_{p}^{*}$ (all nonzero integers modulo $p$).

			      \item Given:
			            \begin{itemize}
				            \item $g \in \mathbb{Z}_{p}^{*}$, $h \in \mathbb{Z}_{p}^{*}$

				            \item Find $x \in \{0, 1, \ldots, p-2\}$ such that $g^{x} \equiv h \pmod
					                  {p}$
			            \end{itemize}

			      \item This problem is believed to be computationally hard when $p$ is large
			            and $g$ is a primitive root modulo $p$.
			            \begin{itemize}
				            \item ``Believed to be'' = we don't know of any way to do it that doesn't
				                  take forever, unless we have a strong, stable quantum computer (Shor's
				                  algorithm)
			            \end{itemize}
		      \end{itemize}
	\end{itemize}
\end{frame}

\begin{frame}{Time for more actual magic}
	\begin{columns}[c]
		\begin{column}{0.6\textwidth}
			\begin{itemize}[<+->]
				\item \textbf{Zero-knowledge proofs} allow you to prove that you know
				      a secret without revealing any information about it.
				\item They built ``zero-knowledge virtual machines'' where you can execute
				      an entire program that runs as a zero-knowledge proof.
				\item ZKP battleship game: server proves to the players that its
				      output to their battleship guesses is correct, without revealing any
				      additional information (e.g. ship location).
			\end{itemize}
		\end{column}

		\begin{column}{0.4\textwidth}
			\imagewithcaption{battleship.jpg}{Battleship board game. Source: Hasbro}
		\end{column}
	\end{columns}
\end{frame}

\begin{frame}{Hard problems}
	\begin{columns}[c]
		\begin{column}{0.5\textwidth}
			\textbf{Asymmetric Primitives}
			\begin{itemize}[<+->]
				\item Diffie-Hellman, RSA, ML-KEM, etc.
				\item ``Asymmetric'' because there is a ``public key'' and a ``private
				      key'' for each party.
				\item Algebraic, assume the hardness of mathematical problems (as seen
				      just now.)
			\end{itemize}
		\end{column}

		\begin{column}{0.5\textwidth}
			\textbf{Symmetric Primitives}
			\begin{itemize}[<+->]
				\item AES, SHA-2, ChaCha20, HMAC\ldots
				\item ``Symmetric'' because there is one secret key.
				\item Not algebraic but unstructured, but on their understood
				      resistance to $n$ years of cryptanalysis.
				\item Can act as substitutes for assumptions in security proofs!
				      \begin{itemize}
					      \item Example: hash function assumed to be a ``random oracle''
				      \end{itemize}
			\end{itemize}
		\end{column}
	\end{columns}
\end{frame}

\begin{frame}{Hard problems}
	\begin{itemize}[<+->]
		\item Hard computational problems are the cornerstone of modern cryptography.
		\item These are problems for which even the best algorithms wouldn't find a solution before the sun burns out.
		\item They provide the security foundation for cryptographic schemes.
		\item Without hard problems, most of our encryption systems would collapse.
	\end{itemize}
\end{frame}

\begin{frame}{The rise of computational complexity theory}
	\definitionbox{Computational Complexity Theory}{Complexity theory provides the mathematical framework to understand what makes problems ``hard''.}
	\begin{itemize}
		\item In the 1970s, rigorous study of hard problems led to computational complexity theory.
		\item This field has had dramatic impacts beyond cryptography:
		      \begin{itemize}
			      \item \textbf{Economics}: Computational complexity of finding Nash equilibria in game theory.
			      \item \textbf{Physics}: Simulating quantum many-body systems with exponential complexity.
			      \item \textbf{Biology}: Protein folding prediction and DNA sequence alignment algorithms.
		      \end{itemize}
	\end{itemize}
\end{frame}

\begin{frame}{Computational problems}
	\definitionbox{Computational Problem}{
		A question that can be answered by performing a computation.
		\begin{itemize}
			\item \textbf{Decision problems}: Questions with ``yes'' or ``no'' answers
			      \begin{itemize}
				      \item Example: ``Is 217 a prime number?''
			      \end{itemize}
			\item \textbf{Search problems}: Questions that require finding a specific value
			      \begin{itemize}
				      \item Example: ``How many instances of \textit{`i'}s appear in \textit{`incomprehensibilities'}?''
			      \end{itemize}
		\end{itemize}
	}
	\begin{itemize}[<+->]
		\item Computational problems form the foundation of theoretical computer science.
		\item Different types of problems require different algorithmic approaches.
		\item The difficulty of solving these problems is central to cryptography.
	\end{itemize}
\end{frame}

\begin{frame}{Computational hardness}
	\definitionbox{Computational Hardness}{
		The property of computational problems for which no algorithm exists that can solve the problem in a reasonable amount of time.
		\begin{itemize}
			\item Also called \textbf{intractable problems}.
			\item Hardness is independent of the computing device used.
			\item All standard computing models are equivalent in terms of what they can compute efficiently.
			\item \textbf{Exception}: Quantum computers for certain problems.
		\end{itemize}
	}
	\begin{itemize}[<+->]
		\item Hardness is a fundamental concept in computational complexity theory.
		\item Cryptography deliberately uses hard problems to create security.
		\item What's ``hard'' should remain hard regardless of hardware advances.
	\end{itemize}
\end{frame}

\begin{frame}{Measuring algorithm complexity}
	\begin{itemize}[<+->]
		\item To evaluate computational hardness, we need to measure an algorithm's running time.
		\item We typically use \textbf{asymptotic analysis} to express complexity
		\item Common notation:
		      \begin{itemize}
			      \item $O(n)$: Linear time.
			      \item $O(n^2)$: Quadratic time.
			      \item $O(2^n)$: Exponential time.
		      \end{itemize}
		\item We care about how the running time grows as the input size increases.
		\item \textbf{Example}: An algorithm that takes $n^2$ operations for input size $n$ becomes impractical as $n$ grows large.
	\end{itemize}
\end{frame}

\begin{frame}{Categorizing computational hardness}
	\begin{columns}[c]
		\begin{column}{0.5\textwidth}
			\textbf{Easy Problems}
			\begin{itemize}[<+->]
				\item Solvable in polynomial time.
				\item \textbf{Examples}: Sorting, searching.
				\item Running time: $O(n^c)$ for some constant $c$
				\item Generally scales reasonably with input size.
				\item Class P (Polynomial time).
			\end{itemize}
		\end{column}
		\begin{column}{0.5\textwidth}
			\textbf{Hard Problems}
			\begin{itemize}[<+->]
				\item No known polynomial-time solution.
				\item \textbf{Example}: Factorizing product of two large primes.
				\item Running time: Often exponential, e.g., $O(2^n)$
				\item Becomes impractical quickly as input grows.
				\item Includes NP-hard, NP-complete classes.
			\end{itemize}
		\end{column}
	\end{columns}
\end{frame}

\begin{frame}{Hard problems in practice}
	\begin{itemize}
		\item Public-key cryptography relies on specific hard problems:
		      \begin{itemize}
			      \item RSA: Integer factorization problem.
			      \item Diffie-Hellman: Discrete logarithm problem.
		      \end{itemize}
		\item Cryptography leverages these problems to maximize security assurance,
		\item The security of these schemes depends on the continued hardness of these problems.
	\end{itemize}
\end{frame}

\begin{frame}{Quantum vulnerability of hard problems}
	\begin{itemize}[<+->]
		\item The hard problems we rely on today (factoring, discrete logarithm) are vulnerable to quantum computers.
		\item Shor's algorithm (1994) can efficiently solve both problems on a sufficiently powerful quantum computer.
		\item This has motivated the search for \textbf{``post-quantum''} hard problems:
		      \begin{itemize}[<+->]
			      \item Lattice-based cryptography (e.g., ML-KEM, formerly CRYSTALS-Kyber).
			      \item Hash-based cryptography.
			      \item Code-based cryptography.
			      \item Multivariate cryptography.
			      \item Isogeny-based cryptography.
		      \end{itemize}
		\item NIST is currently standardizing post-quantum cryptographic algorithms to replace our vulnerable systems.
	\end{itemize}
\end{frame}

\begin{frame}{What is NIST?}
	\begin{columns}[c]
		\begin{column}{0.6\textwidth}
			\begin{itemize}[<+->]
				\item \textbf{NIST} stands for the National Institute of Standards and Technology.
				\item It's a U.S. government agency that develops technology standards.
				\item In cryptography, NIST:
				      \begin{itemize}
					      \item Sets security standards used worldwide.
					      \item Evaluates and approves cryptographic algorithms.
					      \item Currently leading the standardization of post-quantum cryptography.
				      \end{itemize}
				\item When NIST standardizes an algorithm, it often becomes the global industry standard.
			\end{itemize}
		\end{column}
		\begin{column}{0.4\textwidth}
			\imagewithcaption{nist_peanut.png}{NIST's ``Standard Reference Peanut Butter'', availablefor only \$1,217 USD!}
		\end{column}
	\end{columns}
\end{frame}

\begin{frame}{Funny things standardized by NIST}
	\begin{itemize}[<+->]
		\item \textbf{Standard Reference Peanut Butter}: for calibrating food testing equipment.
		\item \textbf{The ``Odor Unit''}: for standardizing measurements of smell intensity in environmental monitoring.
		\item \textbf{The Standard Banana Equivalent Dose (BED)}: for comparing radiation exposure levels to the natural radiation in a banana.
		\item \textbf{Toilet Paper Testing}: for measuring strength, absorbency, and softness of toilet paper products.
	\end{itemize}
\end{frame}

\begin{frame}{Cryptographic algorithms standardized by NIST}
	\begin{itemize}[<+->]
		\item \textbf{AES (Advanced Encryption Standard)}: Selected in 2001 to replace DES, now the worldwide standard for symmetric encryption.
		\item \textbf{SHA-2 and SHA-3 (Secure Hash Algorithms)}: Cryptographic hash functions used for digital signatures and data integrity.
		\item \textbf{DSA and ECDSA}: Digital Signature Algorithms based on the discrete logarithm problem.
		\item \textbf{Triple DES}: An interim standard before AES that enhanced the security of the original DES.
		\item \textbf{ML-KEM and ML-DSA}: Recently standardized post-quantum public-key cryptography and signature schemes.
	\end{itemize}
\end{frame}

\begin{frame}{Why hard problems matter}
	\begin{itemize}
		\item Hard problems \textbf{create asymmetry between legitimate users and attackers}.
		\item Easy in one direction, difficult in the reverse.
		\item Example: Easy to multiply large primes, hard to factor the product.
		\item This asymmetry is what enables secure communication!
	\end{itemize}
\end{frame}

\begin{frame}{Slides not complete}
	\begin{itemize}
		\item This slide deck is not finished and is missing important material. Do not rely on it yet.
	\end{itemize}
\end{frame}

\begin{frame}[plain]
	\titlepage
\end{frame}
\end{document}
